% ============================================================
\chapter{Variables Al\'eatoires Discr\`etes}
\label{ch:discretes}
% ============================================================

Lancer un d\'e, compter le nombre de clients dans une file d'attente, observer le nombre de d\'esint\'egrations radioactives par seconde~: dans chacune de ces situations, le hasard produit un r\'esultat qu'on peut \'enum\'erer. Les variables al\'eatoires discr\`etes --- celles qui ne prennent qu'un nombre fini ou d\'enombrable de valeurs --- sont les premi\`eres que l'on rencontre, et les plus intuitives. Pourtant, elles r\'eservent des surprises~: la loi de Poisson, d\'ecouverte par Sim\'eon Denis Poisson en 1837 pour mod\'eliser les erreurs judiciaires, r\'eappara\^it partout, des centres d'appels aux impacts de bombes V-2 sur Londres. La loi g\'eom\'etrique capture la patience du joueur obstin\'e, tandis que la loi binomiale, h\'eriti\`ere directe des travaux de Jacob Bernoulli, reste le mod\`ele de r\'ef\'erence pour les exp\'eriences r\'ep\'et\'ees.

Ce chapitre construit le cadre formel des variables al\'eatoires discr\`etes et d\'eveloppe les lois classiques qui peuplent toute la th\'eorie des probabilit\'es.

\begin{intuition}
Une variable al\'eatoire est une fonction qui associe une valeur num\'erique au
r\'esultat d'une exp\'erience al\'eatoire. Ce chapitre traite le cas discret~:
l'ensemble des valeurs prises est fini ou d\'enombrable.
\end{intuition}

% ------------------------------------------------------------
\section{D\'efinition et loi d'une variable al\'eatoire}
% ------------------------------------------------------------

\begin{definition}[Variable al\'eatoire]
Soit $(\Omega, \mathcal{F}, \PP)$ un espace probabilis\'e. Une \textbf{variable
al\'eatoire r\'eelle} est une application $X : \Omega \to \R$ mesurable, c'est-\`a-dire
telle que pour tout $B \in \mathcal{B}(\R)$~:
\[
    X^{-1}(B) = \{\omega \in \Omega : X(\omega) \in B\} \in \mathcal{F}.
\]
\end{definition}

\begin{definition}[Variable al\'eatoire discr\`ete]
Une variable al\'eatoire $X$ est \textbf{discr\`ete} si elle prend ses valeurs dans
un ensemble fini ou d\'enombrable $\{x_1, x_2, \ldots\} \subseteq \R$.
\end{definition}

\begin{definition}[Loi de probabilit\'e / fonction de masse]
La \textbf{loi} (ou \textbf{distribution}) de $X$ est la mesure de probabilit\'e
$P_X$ sur $(\R, \mathcal{B}(\R))$ d\'efinie par $P_X(B) = \PP(X \in B)$.
Pour une variable discr\`ete, la loi est enti\`erement d\'etermin\'ee par la
\textbf{fonction de masse}~:
\[
    p_X(x_k) = \PP(X = x_k), \quad k = 1, 2, \ldots
\]
avec $p_X(x_k) \geq 0$ et $\sum_k p_X(x_k) = 1$.
\end{definition}

\begin{definition}[Fonction de r\'epartition]
La \textbf{fonction de r\'epartition} de $X$ est~:
\[
    F_X(x) = \PP(X \leq x) = \sum_{x_k \leq x} p_X(x_k), \quad x \in \R.
\]
\end{definition}

\begin{proposition}[Propri\'et\'es de $F_X$]
\begin{enumerate}[label=(\roman*)]
    \item $F_X$ est croissante.
    \item $\lim_{x \to -\infty} F_X(x) = 0$ et $\lim_{x \to +\infty} F_X(x) = 1$.
    \item $F_X$ est continue \`a droite.
    \item $\PP(X = a) = F_X(a) - F_X(a^-)$ (saut en $a$).
    \item Pour une variable discr\`ete, $F_X$ est une fonction en escalier.
\end{enumerate}
\end{proposition}

% ------------------------------------------------------------
\section{Lois discr\`etes classiques}
% ------------------------------------------------------------

\subsection{Loi de Bernoulli}

\begin{definition}[Bernoulli]
$X \sim \text{Bernoulli}(p)$ avec $p \in [0,1]$~:
\[
    \PP(X = 1) = p, \quad \PP(X = 0) = 1 - p = q.
\]
$\E[X] = p$, $\Var(X) = pq$.
\end{definition}

\subsection{Loi binomiale}

\begin{definition}[Binomiale]
$X \sim \text{Bin}(n, p)$~: nombre de succ\`es en $n$ essais de Bernoulli
ind\'ependants.
\[
    \PP(X = k) = \binom{n}{k} p^k (1-p)^{n-k}, \quad k = 0, 1, \ldots, n.
\]
$\E[X] = np$, $\Var(X) = np(1-p)$.
\end{definition}

\begin{exemple}
On lance 10 fois une pi\`ece \'equilibr\'ee. La probabilit\'e d'obtenir exactement
6~Pile est~:
\[
    \PP(X = 6) = \binom{10}{6} \left(\frac{1}{2}\right)^{10}
    = \frac{210}{1024} \approx 0{,}205.
\]
\end{exemple}

\subsection{Loi g\'eom\'etrique}

\begin{definition}[G\'eom\'etrique]
$X \sim \text{Geom}(p)$~: nombre d'essais n\'ecessaires pour obtenir le premier
succ\`es.
\[
    \PP(X = k) = (1-p)^{k-1} p, \quad k = 1, 2, 3, \ldots
\]
$\E[X] = 1/p$, $\Var(X) = (1-p)/p^2$.
\end{definition}

\begin{proposition}[Perte de m\'emoire]
La loi g\'eom\'etrique est la seule loi discr\`ete \`a support dans $\N^*$
v\'erifiant la propri\'et\'e de \textbf{perte de m\'emoire}~:
\[
    \PP(X > m + n \mid X > m) = \PP(X > n) \quad \forall m, n \geq 0.
\]
\end{proposition}

\begin{proof}
$\PP(X > k) = (1-p)^k$. Donc~:
\[
    \PP(X > m+n \mid X > m) = \frac{(1-p)^{m+n}}{(1-p)^m} = (1-p)^n = \PP(X > n). \qedhere
\]
\end{proof}

\subsection{Loi de Poisson}

\begin{definition}[Poisson]
$X \sim \text{Poisson}(\lambda)$ avec $\lambda > 0$~:
\[
    \PP(X = k) = e^{-\lambda} \frac{\lambda^k}{k!}, \quad k = 0, 1, 2, \ldots
\]
$\E[X] = \lambda$, $\Var(X) = \lambda$.
\end{definition}

\begin{theoreme}[Approximation de Poisson]
Si $X_n \sim \text{Bin}(n, p_n)$ avec $n p_n \to \lambda$ quand $n \to \infty$,
alors pour tout $k \in \N$~:
\[
    \PP(X_n = k) \xrightarrow[n \to \infty]{} e^{-\lambda} \frac{\lambda^k}{k!}.
\]
\end{theoreme}

\begin{proof}
\begin{align*}
    \PP(X_n = k) &= \binom{n}{k} p_n^k (1-p_n)^{n-k} \\
    &= \frac{n!}{k!(n-k)!} \left(\frac{\lambda}{n}\right)^k
       \left(1 - \frac{\lambda}{n}\right)^{n-k} \\
    &= \frac{\lambda^k}{k!} \cdot \underbrace{\frac{n(n-1)\cdots(n-k+1)}{n^k}}_{\to 1}
       \cdot \underbrace{\left(1 - \frac{\lambda}{n}\right)^n}_{\to e^{-\lambda}}
       \cdot \underbrace{\left(1 - \frac{\lambda}{n}\right)^{-k}}_{\to 1}
    \xrightarrow[n \to \infty]{} e^{-\lambda} \frac{\lambda^k}{k!}. \qedhere
\end{align*}
\end{proof}

\begin{intuition}[title=\textbf{Quand utiliser Poisson~?}]
La loi de Poisson mod\'elise le nombre d'\'ev\'enements rares dans un intervalle
donn\'e (appels t\'el\'ephoniques, particules radioactives, pannes\ldots). R\`egle
pratique~: si $n \geq 30$ et $p \leq 0{,}1$, on approxime $\text{Bin}(n,p)$ par
$\text{Poisson}(np)$.
\end{intuition}

\subsection{Loi hyperg\'eom\'etrique}

\begin{definition}[Hyperg\'eom\'etrique]
On tire $n$ objets sans remise d'une population de $N$ objets dont $K$ sont
marqu\'es. $X$ = nombre d'objets marqu\'es parmi les $n$ tir\'es~:
\[
    \PP(X = k) = \frac{\binom{K}{k}\binom{N-K}{n-k}}{\binom{N}{n}},
    \quad k = \max(0, n+K-N), \ldots, \min(n, K).
\]
$\E[X] = nK/N$.
\end{definition}

% ------------------------------------------------------------
\section{Tableau r\'ecapitulatif}
% ------------------------------------------------------------

\begin{formules}[title=\textbf{Lois discr\`etes fondamentales}]
\begin{center}
\renewcommand{\arraystretch}{1.4}
\begin{tabular}{lccc}
    \toprule
    \textbf{Loi} & \textbf{Param\`etres} & $\E[X]$ & $\Var(X)$ \\
    \midrule
    Bernoulli$(p)$ & $p \in [0,1]$ & $p$ & $p(1-p)$ \\
    Binomiale$(n,p)$ & $n \in \N^*,\; p \in [0,1]$ & $np$ & $np(1-p)$ \\
    G\'eom\'etrique$(p)$ & $p \in (0,1]$ & $1/p$ & $(1-p)/p^2$ \\
    Poisson$(\lambda)$ & $\lambda > 0$ & $\lambda$ & $\lambda$ \\
    Uniforme$\{1,\ldots,n\}$ & $n \in \N^*$ & $(n+1)/2$ & $(n^2-1)/12$ \\
    \bottomrule
\end{tabular}
\end{center}
\end{formules}

% ------------------------------------------------------------
\section{Fonction g\'en\'eratrice}
% ------------------------------------------------------------

\begin{definition}[Fonction g\'en\'eratrice]
Soit $X$ une variable al\'eatoire \`a valeurs dans $\N$. La \textbf{fonction
g\'en\'eratrice} de $X$ est~:
\[
    G_X(s) = \E[s^X] = \sum_{k=0}^{\infty} \PP(X = k)\, s^k, \quad \abs{s} \leq 1.
\]
\end{definition}

\begin{proposition}[Propri\'et\'es de $G_X$]
\begin{enumerate}[label=(\roman*)]
    \item $G_X(1) = 1$.
    \item $\E[X] = G_X'(1)$.
    \item $\E[X(X-1)] = G_X''(1)$, d'o\`u
          $\Var(X) = G_X''(1) + G_X'(1) - [G_X'(1)]^2$.
    \item $\PP(X = k) = G_X^{(k)}(0)/k!$.
    \item Si $X$ et $Y$ sont ind\'ependantes, $G_{X+Y}(s) = G_X(s) \cdot G_Y(s)$.
\end{enumerate}
\end{proposition}

\begin{exemple}[Fonction g\'en\'eratrice de la loi de Poisson]
Si $X \sim \text{Poisson}(\lambda)$~:
\[
    G_X(s) = \sum_{k=0}^{\infty} e^{-\lambda} \frac{\lambda^k}{k!} s^k
    = e^{-\lambda} \sum_{k=0}^{\infty} \frac{(\lambda s)^k}{k!}
    = e^{-\lambda} \cdot e^{\lambda s}
    = e^{\lambda(s-1)}.
\]
On retrouve $G_X'(1) = \lambda = \E[X]$ et
$G_X''(1) = \lambda^2$, d'o\`u $\Var(X) = \lambda^2 + \lambda - \lambda^2 = \lambda$.
\end{exemple}

% ------------------------------------------------------------
\section{Exercices}
% ------------------------------------------------------------

\begin{exercice}
Soit $X \sim \text{Bin}(n, p)$. Montrer directement (par calcul) que
$\E[X] = np$ et $\Var(X) = np(1-p)$.
\end{exercice}

\begin{exercice}
Soit $X \sim \text{Geom}(p)$. Calculer $\E[X]$ en utilisant la fonction
g\'en\'eratrice.
\end{exercice}

\begin{exercice}
Soit $X \sim \text{Poisson}(\lambda)$ et $Y \sim \text{Poisson}(\mu)$,
ind\'ependantes. Montrer que $X + Y \sim \text{Poisson}(\lambda + \mu)$ en
utilisant les fonctions g\'en\'eratrices.
\end{exercice}

\begin{exercice}
Un standard t\'el\'ephonique re\c{c}oit en moyenne 4~appels par minute, mod\'elis\'es
par une loi de Poisson. Calculer la probabilit\'e de recevoir~:
\begin{enumerate}[label=(\alph*)]
    \item exactement 3~appels en une minute\,;
    \item au moins 1~appel en 30~secondes\,;
    \item plus de 10~appels en 2~minutes.
\end{enumerate}
\end{exercice}

\begin{exercice}
Un sac contient 5~boules rouges et 15~boules bleues. On tire 4~boules sans remise.
Calculer la probabilit\'e d'obtenir exactement 2~boules rouges.
\end{exercice}

\begin{exercice}
On lance un d\'e \'equilibr\'e jusqu'\`a obtenir un~6. Soit $X$ le nombre de lancers.
\begin{enumerate}[label=(\alph*)]
    \item Quelle est la loi de $X$~?
    \item Calculer $\PP(X > 10)$.
    \item Calculer $\PP(X > 15 \mid X > 5)$. Retrouver la perte de m\'emoire.
\end{enumerate}
\end{exercice}

\begin{exercice}
Montrer que la loi g\'eom\'etrique est la seule loi discr\`ete sur $\N^*$ v\'erifiant
la propri\'et\'e de perte de m\'emoire.
\textit{Indication~:} poser $g(n) = \PP(X > n)$ et utiliser l'\'equation fonctionnelle
$g(m+n) = g(m)g(n)$.
\end{exercice}

\begin{exercice}
Soit $N \sim \text{Poisson}(\lambda)$ et $(X_i)_{i \geq 1}$ des variables i.i.d.\
de Bernoulli$(p)$, ind\'ependantes de $N$. Poser $S = \sum_{i=1}^{N} X_i$.
Montrer que $S \sim \text{Poisson}(\lambda p)$.
\end{exercice}
